%% !TEX root = ../manual_gaolv5.tex
% Chapter 10.
%% Chapter 10: The names of IEEE 1788-2015. Sources: gaol.tex of f5bf180 lines
%% 5532-5788 (the chapter "The names of IEEE 1788-2015" of the current
%% manual), doc/using.md lines 451-571 (the names of IEEE 1788-2015, the
%% reverse functions of atan2).
%% Based on the manual of GAOL 4 by Frederic Goualard (gaol.tex).
%%
%% Copyright (c) 2026 ENSTA, France
%%
%% By: Jordan NININ (created 2026-11-06)
%%
%% GAOL is a software distributed WITHOUT ANY WARRANTY. Read the associated
%% COPYING.LIB file for information.
\chapter{The names of IEEE 1788-2015}\label{chap:ieee1788}
%%===================================
\cindex{namespace!gaol\_ieee1788}\cindex{IEEE 1788-2015!names}%

IEEE 1788-2015 \cite{IEEE1788:2015} names the operations of interval
arithmetic, and GAOL names its operations its own way. The reverse function
\ieee{coshRev}$(c,x)$ of the standard is \code{acosh\_rel(c, x)} in GAOL,
\ieee{mulRev}$(b,c,x)$ is \code{div\_rel(c, b, x)}, with its arguments in
another order, and \ieee{roundTiesToEven}$(x)$ is
\code{round\_ties\_to\_even(x)}. A reader of the standard who comes to GAOL
has to look for the translation of each name.

\medskip
The namespace \code{gaol\_ieee1788} spares this search. It gives each
operation of the standard that GAOL provides the name and the argument order
the standard gives it, so that a reader of the standard finds it without
looking for its translation. It is the second of the two namespaces of GAOL
v5 (\pxref{sec:namespaces}); the standard itself is presented in
Chapter~\ref{chap:concepts} (\pxref{sec:ieee1788-overview}).

\medskip
In this chapter, we describe this namespace. The first section presents it
and the way a program opens it. The second lists the operations whose
meaning differs between the standard and GAOL, and which follow the standard
in \code{gaol\_ieee1788}. The third gives the table of the correspondence
between the names of the standard and those of GAOL, then details
\ieee{mulRevToPair}, which has no name in \code{gaol}, and the reverse
functions of \ieee{atan2}. The fourth describes the reading of intervals
with the names of the standard, and the fifth lists what GAOL does not
provide of the standard. The names of the other interval libraries, next to
those of GAOL and of the standard, are compared in
Chapter~\ref{chap:overview} (\pxref{sec:other-libraries}).

\section{The namespace of the standard}\label{sec:ieee1788-namespace}
%%------------------------------------

The namespace \code{gaol\_ieee1788} comes with \file{gaol/gaol.h}. It is not in
\code{gaol}, and it holds
the type \code{interval} as well: one directive lets a program use GAOL under
the names of the standard. The output format is not in it: a program that opens
only \code{gaol\_ieee1788} sets it with
\code{interval::format(gaol::interval\_format::hexa)}.\newinvfive

\begin{Ex}\label{ex:ieee1788-program}
The program below opens \code{gaol\_ieee1788} and computes with the names of
the standard: \code{numsToInterval} builds an interval from its two bounds,
\code{mulRev} is the reverse multiplication, which divides \code{x} by
\itv{2}{2}, \code{sinPi} and \code{rootn} are two functions of Table~10.5 of
the standard, \code{entire} is the interval \itv{-\infty}{+\infty} and
\code{strictLess} the strict order of its Table~10.3. What it prints
follows it.
\begin{example}
#include <iostream>
#include <gaol/gaol.h>
using namespace gaol_ieee1788;

int main(void)
{
  interval x = numsToInterval(1, 2);
  interval y = mulRev(numsToInterval(2, 2), x);   @rem^// x / 2~
  interval z = sinPi(x) + rootn(y, 3);
  bool b = strictLess(x, entire());
  std::cout << std::boolalpha
            << x << ' ' << y << ' ' << z << ' ' << b << std::endl;
  gaol::cleanup();
  return 0;
}
@outputs^[1, 2] [0.5, 1] [-0.2062994740159004, 1] false~
\end{example}
\end{Ex}

\medskip
\noindent
A program opens \code{gaol} or \code{gaol\_ieee1788}, not both
(\pxref{sec:namespaces}). The functions that have the same name and the same
meaning in both are the same functions; the ones whose meaning differs,
\code{pow} and \code{textToInterval}, exist in each namespace, and a program
that opened both would find each call ambiguous.

\begin{Rem}\label{rem:ieee1788-shadow}
A name of the program's own that is also a name of the standard, a constant
\code{inf} for instance, makes every unqualified use of it ambiguous: the
program qualifies both, \code{::inf} and \code{gaol\_ieee1788::inf(x)}.
The same holds for \code{less(x, y)} next to \code{using namespace std;}, where it is
ambiguous with the class template \code{std::less}.
\end{Rem}

\section{Where the standard and GAOL differ}\label{sec:ieee1788-differences}
%%-----------------------------------------
\cindex{IEEE 1788-2015!differences with GAOL}%

The functions of GAOL that have the name and the meaning the standard gives
them (\code{sin}, \code{exp}, \code{sqrt}, \code{min}\ldots) are the same
functions in \code{gaol\_ieee1788}. Where the standard and GAOL differ, the
names of \code{gaol\_ieee1788} follow the standard. The differences are the
following ones.

\begin{itemize}
\item[$\bullet$] \textbf{The power.} \code{pow(x, y)} is the \ieee{pow} of
  Table~9.1, defined for $x>0$, and for $x=0$ when $y>0$. The \code{pow} of
  \code{gaol} takes the integer power for an integer exponent, a negative
  base included (\pxref{sec:powers}). \code{pow(x, p)} with a number $p$ is
  \code{pow(x, [p])}, an integer $p$ included, \code{pow(x, 2)} as well, as
  the standard gives no other meaning to it: the power with an integer
  exponent is \code{pown(x, n)} in the standard, whose $n$ is an integer
  rather than an interval. An integer exponent beyond the \code{int}s gives
  the power of CORE-MATH at the bounds of the part of $x$ in
  \itv{0}{+\infty}, as the \code{pow} of \code{gaol} does there; the \code{pow}
  of \code{gaol} takes the integer power on the negative part of $x$ too. The results of the two powers are
  compared in Chapter~\ref{chap:overview} (\pxref{sec:pow-namespaces}).
\item[$\bullet$] \textbf{The bounds of the empty set.} \code{inf(x)} and
  \code{sup(x)} are $+\infty$ and $-\infty$ for the empty set (Table~10.2),
  where \code{left()} and \code{right()} are NaN. \code{inf} returns $-0$
  for a left bound 0 (12.12.8).
\item[$\bullet$] \textbf{Membership.} \code{isMember(m, x)} is \code{false}
  for an infinite $m$ (10.6.3).
\item[$\bullet$] \textbf{Writing.} \code{intervalToText(x)} writes the
  bounds of \code{x} rounded outward, \code{[l, r]}, or \code{[a]} for a
  degenerate interval whose value the digits write exactly, as \code{[4]},
  or \code{[empty]}, with the digits of \code{interval::precision()} and a
  decimal point. This string is an interval literal of the standard (12.11),
  which \code{textToInterval} reads back as an interval containing
  \code{x} (13.3), whatever the output format that
  \code{interval::format()} sets and the locale of the program.
  \code{operator<{}<} writes the current format instead, which is no
  literal in the formats \element{width} and \element{center}, nor in the
  format \element{agreeing} where the bounds share digits.\newinvfive
\item[$\bullet$] \textbf{Reading.} \code{textToInterval} and
  \code{exactToInterval} read the names of the functions of the standard,
  and return the empty set for a string that is no interval (12.1.3), where
  \code{gaol::textToInterval} reads the names of GAOL and throws
  (\pxref{sec:ieee1788-text}). A null pointer is no string, and they throw
  \code{std::invalid\_argument} for it too (\pxref{sec:constructors}).%
  \newinvfive
\end{itemize}

\begin{Ex}\label{ex:ieee1788-differences}
The program below shows these differences. The \code{pow} of the standard is
empty on a negative base, where \code{pown} gives the integer power; an
integer exponent beyond the \code{int}s gives the power of CORE-MATH at the
bounds; the bounds of the empty set are the infinities; an infinity is a
member of no interval; and a string that is no interval reads as the empty
set.
\begin{example}
cout << pow(interval(-4,-1), interval(2.0)) << ' '
     << pown(interval(-4,-1), 2) << ' '
     << pow(interval(2,3), interval(1e10)) << endl;
@outputs^[empty] [1, 16] [1.797693134862315e+308, inf]~
cout << boolalpha << inf(empty()) << ' ' << sup(empty()) << ' '
     << isMember(GAOL_INFINITY, entire()) << ' '
     << textToInterval("[1, 2") << endl;
@outputs^inf -inf false [empty]~
\end{example}
\end{Ex}

\section{The operations}\label{sec:ieee1788-operations}
%%----------------------
\cindex{IEEE 1788-2015!operations}%

Table~\ref{tab:ieee1788-operations} gives, for each operation of the standard
that GAOL provides, its name in \code{gaol\_ieee1788} and the name GAOL gives
it, grouped as the standard groups them, with the clause or the table of the
standard that defines them. The last argument of the reverse functions
is optional, and \itv{-\infty}{+\infty} when it is left out (10.5.4).

\begingroup\renewcommand{\arraystretch}{1.1}\small
\begin{longtable}{L{6.2cm}L{7.8cm}}
  \caption{The operations of IEEE 1788-2015 in \code{gaol\_ieee1788}, and
    their names in GAOL}\label{tab:ieee1788-operations}\\
  \hline
  \multicolumn{1}{c}{\code{gaol\_ieee1788}} & \multicolumn{1}{c}{GAOL}\\
  \hline\hline
  \endfirsthead
  \multicolumn{2}{l}{\textit{Table~\thetable, continued}}\\*
  \hline
  \multicolumn{1}{c}{\code{gaol\_ieee1788}} & \multicolumn{1}{c}{GAOL}\\
  \hline\hline
  \endhead
  \hline
  \endfoot
  \multicolumn{2}{l}{\textit{Constants and constructors (10.5.2, 10.5.8)}}\\*
  \code{empty()}, \code{entire()} &
    \code{interval::emptyset()}, \code{interval::universe()}\\
  \code{numsToInterval(l, u)} & \code{interval(l, u)}\\
  \code{textToInterval(s)} &
    \code{gaol::textToInterval(s)} with the names of the standard, the
    empty set for a wrong string\\
  \hline
  \multicolumn{2}{l}{\textit{Forward elementary functions (Table~9.1 of the standard)}}\\*
  \code{neg}, \code{add}, \code{sub}, \code{mul}, \code{div} &
    \code{-x}, \code{x + y}, \code{x - y}, \code{x * y}, \code{x / y}\\
  \code{recip(x)} & \code{inverse(x)}\\
  \code{sqr}, \code{sqrt}, \code{fma} & the same\\
  \code{pown(x, p)} &
    \code{gaol::pow(x, p)} for an integer $p$ of any type\\
  \code{pow(x, y)}, \code{pow(x, p)} & the \ieee{pow} of the standard\\
  \code{exp}, \code{exp2}, \code{exp10}, \code{log}, \code{log2},
    \code{log10} & the same\\
  \code{sin}, \code{cos}, \code{tan}, \code{asin}, \code{acos},
    \code{atan}, \code{atan2} & the same\\
  \code{sinh}, \code{cosh}, \code{tanh}, \code{asinh}, \code{acosh},
    \code{atanh} & the same\\
  \code{sign}, \code{ceil}, \code{floor}, \code{trunc} & the same\\
  \code{roundTiesToEven}, \code{roundTiesToAway} &
    \code{round\_ties\_to\_even}, \code{round\_ties\_to\_away}\\
  \code{abs}, \code{min}, \code{max} & the same\\
  \hline
  \multicolumn{2}{l}{\textit{Recommended functions (Table~10.5 of the standard)}}\\*
  \code{rootn(x, q)} & \code{nth\_root(x, q)}\\
  \code{expm1}, \code{exp2m1}, \code{exp10m1} & the same\\
  \code{logp1}, \code{log2p1}, \code{log10p1} &
    \code{log1p}, \code{log2p1}, \code{log10p1}\\
  \code{hypot}, \code{rSqrt} & \code{hypot}, \code{rsqrt}\\
  \code{compoundm1} & the same\\
  \code{sinPi}, \code{cosPi}, \code{tanPi} &
    \code{sinpi}, \code{cospi}, \code{tanpi}\\
  \code{asinPi}, \code{acosPi}, \code{atanPi}, \code{atan2Pi} &
    \code{asinpi}, \code{acospi}, \code{atanpi}, \code{atan2pi}\\
  \hline
  \multicolumn{2}{l}{\textit{Reverse functions (Table~10.1 of the standard)}}\\*
  \code{sqrRev(c, x)}, \code{absRev(c, x)} &
    \code{sqrt\_rel(c, x)}, \code{invabs\_rel(c, x)}\\
  \code{pownRev(c, x, p)} & \code{nth\_root\_rel(c, p, x)}, for every integer $p$,
    $p\leq0$ included\\
  \code{sinRev}, \code{cosRev}, \code{tanRev(c, x)} &
    \code{asin\_rel}, \code{acos\_rel}, \code{atan\_rel(c, x)}\\
  \code{coshRev(c, x)} & \code{acosh\_rel(c, x)}\\
  \code{sinhRev}, \code{tanhRev(c, x)} (not in the standard, named after
    \ieee{coshRev}) & \code{asinh\_rel}, \code{atanh\_rel(c, x)}\\
  \code{mulRev(b, c, x)} & \code{div\_rel(c, b, x)}\\
  \code{powRev1(b, c, x)}, \code{powRev2(a, c, x)} &
    \code{pow\_rel(c, b, x)},
    \code{pow\_exponent\_rel(c, a, x)}\newinvfive\\
  \code{atan2Rev1(x, c, y)}, \code{atan2Rev2(y, c, x)} &
    \code{atan2\_rel\_y(c, x, y)},
    \code{atan2\_rel\_x(c, y, x)}\newinvfive\\
  \code{mulRevToPair(b, c)} (10.5.5) &
    the parts of $c/b$ on either side of 0, each a
    \code{div\_rel} with \code{interval::universe()} for its last argument, in a
    \code{std::pair}\\
  \hline
  \multicolumn{2}{l}{\textit{Cancellative operations and set operations
    (10.5.6, 10.5.7)}}\\*
  \code{cancelMinus}, \code{cancelPlus} &
    \code{cancel\_minus}, \code{cancel\_plus}\\
  \code{intersection(x, y)}, \code{convexHull(x, y)} &
    \code{x \& y}, \code{x | y}\\
  \hline
  \multicolumn{2}{l}{\textit{Numeric functions (Table~10.2 of the standard)}}\\*
  \code{inf(x)}, \code{sup(x)} &
    \code{x.left()}, \code{x.right()}, but for the empty set and a zero bound
    ($-0$ for \code{inf}, $+0$ for \code{sup})\\
  \code{mid}, \code{wid}, \code{rad}, \code{mag}, \code{mig(x)} &
    \code{x.midpoint()}, \code{width()}, \code{rad()}, \code{mag()},
    \code{mig()}\\
  \code{midRad(x, m, r)} & \code{x.mid\_rad(m, r)}\\
  \hline
  \multicolumn{2}{l}{\textit{Boolean functions (10.5.10, Tables~10.3 and
    10.4 of the standard; 10.6.3)}}\\*
  \code{isEmpty(x)}, \code{isEntire(x)} &
    \code{x.is\_empty()}, \code{x.is\_entire()}\\
  \code{equal(a, b)}, \code{subset(a, b)} &
    \code{a.set\_eq(b)}, \code{a.set\_leq(b)}\\
  \code{interior(a, b)}, \code{disjoint(a, b)} &
    \code{a.set\_le(b)}, \code{a.set\_disjoint(b)}\\
  \code{less(a, b)}, \code{strictLess(a, b)} &
    \code{a.less(b)}, \code{a.strictly\_less(b)}\\
  \code{precedes(a, b)}, \code{strictPrecedes(a, b)} &
    \code{a.certainly\_leq(b)}, \code{a.certainly\_le(b)}\\
  \code{isCommonInterval(x)}, \code{isSingleton(x)} &
    \code{x.is\_common\_interval()}, \code{x.is\_a\_double()}\\
  \code{isMember(m, x)} &
    \code{x.set\_contains(m)}, false for an infinite $m$\\
  \code{overlap(a, b)} (10.6.4) &
    \code{a.overlap(b)}, an \code{overlapping\_state}\\
  \hline
  \multicolumn{2}{l}{\textit{Input and output (Clause 13)}}\\*
  \code{intervalToText(x)} &
    the bounds \code{[l, r]}, \code{[a]} or \code{[empty]}, with the
    current precision, whatever the format and locale\\
  \code{intervalToExact(x)} & \code{exact\_string(x)}\\
  \code{exactToInterval(s)} & \code{textToInterval(s)} of \code{gaol\_ieee1788}\\
  \hline
  \multicolumn{2}{l}{\textit{Expressions}}\\*
  \code{pown(e, n)}, \code{pow(e1, e2)} &
    the expressions of \code{pown} and of the \ieee{pow} of the standard,
    declared by \file{gaol/parser/\allowbreak gaol\_expression.h}\\
\end{longtable}
\endgroup

\noindent
Two rows of the table call for a word. \code{overlap(a, b)} is the interval
overlapping function of the standard: it returns an
\code{overlapping\_state}, one of the sixteen states of Table~10.7 under the
names of the table, and the type and the member \code{a.overlap(b)} are in
\code{gaol} too, which has no other name for it
(Chapter~\ref{chap:relations}).\newinvfive{} The expressions \code{pown(e, n)}
and \code{pow(e1, e2)} are computed with the \code{pown} and the \code{pow}
of this namespace; \file{gaol/parser/gaol\_expression.h} declares them with
the other expressions, and \file{gaol/gaol.h} does not include it, as in
GAOL~4. The last two parts of this section detail \code{mulRevToPair} and
the reverse functions of \code{atan2}.

\paragraph{The two-output division.}
\findex{mulRevToPair}%
\code{mulRevToPair(b, c)} returns the set
$\{x \mid \exists\beta\in b,\ \beta x\in c\}$ in two intervals, the pair
$(u,v)$ of the standard: $u$ before $v$ when the set has two parts, which
happens when $c$ does not contain 0 and 0 is strictly between the bounds of
$b$, $(u,\emptyset)$ when it is one interval, and $(\emptyset,\emptyset)$ when
it is empty. The convex hull of $u$ and $v$ is \code{mulRev(b, c)}. Where
$b$ is unbounded on both sides, the two parts share 0.\newinvfive

\begin{Ex}\label{ex:ieee1788-mulrevtopair}
The program below divides \itv{1}{2} by \itv{-1}{1}, whose interior holds 0,
and gets two parts; then by \itv{2}{4}, and gets one part and the empty set.
\begin{example}
std::pair<interval, interval> p = mulRevToPair(numsToInterval(-1, 1),
                                               numsToInterval(1, 2));
std::pair<interval, interval> q = mulRevToPair(numsToInterval(2, 4),
                                               numsToInterval(1, 2));
cout << p.first << ' ' << p.second << ' '
     << q.first << ' ' << q.second << endl;
@outputs^[-inf, -1] [1, inf] [0.25, 1] [empty]~
\end{example}
\end{Ex}

\paragraph{The reverse functions of atan2.}
\findex{atan2Rev1}\findex{atan2Rev2}%
\code{atan2Rev1(x, c, y)} is the hull of the ordinates of the points of the
box $x\times y$ whose angle $\operatorname{atan2}(y,x)$ is in $c$, and
\code{atan2Rev2(y, c, x)} the hull of their abscissas; the point $(0,0)$,
which has no angle, is left out. In \code{gaol}, they are
\code{atan2\_rel\_y(c, x, y)} and \code{atan2\_rel\_x(c, y, x)}
(\pxref{sec:relational-arithmetic}).\newinvfive{} As in \code{mulRev} and
\code{div\_rel}, the interval contracted comes last, and the other argument
of \code{atan2} before it, so that the constraint
$\theta=\operatorname{atan2}(y,x)$ contracts its three variables with the
three lines below.
\begin{example}
theta &= atan2(y, x);
y = atan2Rev1(x, theta, y);
x = atan2Rev2(y, theta, x);
\end{example}

\noindent
The angles are those of \code{atan2}, in $\left]-\pi,\pi\right]$. The
half-line $y=0$, $x<0$ has the angle $\pi$, not $-\pi$, and a $c$ holding
$-\pi$ but not $\pi$ leaves it out: \code{atan2Rev2([0], [-4, -3], [-2, 2])}
is empty. The angles of $c$ beyond $\pi$ are no angle of any point:
\code{atan2Rev1([-1], [3, 4], [-1, 1])} is \itv{0}{-\tan 3}, about
\itv{0}{0.1425}, the points of angles 3 to $\pi$ only. A program whose
angles go around past $\pi$ gives the two sides of the jump apart,
\itv{3}{\pi} and \itv{-\pi}{4-2\pi} for \itv{3}{4}, and takes the hull of
the two results.

\section{Reading with the names of the standard}\label{sec:ieee1788-text}
%%-----------------------------------------------
\findex{textToInterval}\cindex{IEEE 1788-2015!reading}%

\code{textToInterval(\var{s})} reads a string as \code{gaol::textToInterval()}
does (\pxref{sec:input-format}), with the names of the functions of IEEE
1788-2015: those of Tables~9.1 and 10.5 that GAOL provides, whatever the case
of their letters.\newinvfive{} These names are the following ones.

\begin{center}\small
\begin{tabular}{L{11cm}}
\code{neg}, \code{add}, \code{sub}, \code{mul}, \code{div}, \code{recip},
\code{sqr}, \code{sqrt}, \code{fma}, \code{pown}, \code{pow}, \code{exp},
\code{exp2}, \code{exp10}, \code{log}, \code{log2}, \code{log10}, \code{sin},
\code{cos}, \code{tan}, \code{asin}, \code{acos}, \code{atan}, \code{atan2},
\code{sinh}, \code{cosh}, \code{tanh}, \code{asinh}, \code{acosh},
\code{atanh}, \code{sign}, \code{ceil}, \code{floor}, \code{trunc},
\code{roundTiesToEven}, \code{roundTiesToAway}, \code{abs}, \code{min},
\code{max}, \code{rootn}, \code{expm1}, \code{exp2m1}, \code{exp10m1},
\code{logp1}, \code{log2p1}, \code{log10p1}, \code{hypot},
\code{compoundm1}, \code{rSqrt}, \code{sinPi}, \code{cosPi}, \code{tanPi},
\code{asinPi}, \code{acosPi}, \code{atanPi}, \code{atan2Pi}\\
\end{tabular}
\end{center}

\noindent
\code{pow} is the \ieee{pow} of Table~9.1, and the second argument of
\code{pown} and \code{rootn} is an integer. The names of GAOL alone
(\code{nth\_root}, \code{cbrt}, \code{inverse}, \code{integer},
\code{log1p}, \code{round\_ties\_to\_even}\ldots) are no functions there.
They give the empty set, as any string that is no interval: a wrong call
(\code{"pown([2,5],2.5)"}, \code{"sin(1,2)"}) and a string holding a NUL
character included, whether GAOL was configured with exceptions or
not.\newinvfive{} Where memory runs out, \code{textToInterval} throws
\code{std::bad\_alloc}.

\medskip
As \code{pow}, the function \code{textToInterval} is in the namespace
\code{gaol} too, where it reads the names of GAOL and throws
\code{std::invalid\_argument} for a wrong string, and \code{std::domain\_error}
for an order of \code{nth\_root} that is no integer: a program calls the one of
the namespace it opens.

\begin{Ex}\label{ex:ieee1788-text}
The program below reads three strings with the names of the standard: a
power of the standard, a power of a negative base, which the \ieee{pow} of
the standard leaves empty, and a name of GAOL alone, which is no function
there.
\begin{example}
cout << textToInterval("pown([2,5],5)") << ' '
     << textToInterval("pow([-4,-1],2)") << ' '
     << textToInterval("nth_root(8,3)") << endl;
@outputs^[32, 3125] [empty] [empty]~
\end{example}
With \code{using namespace gaol;}, \code{textToInterval("pow([-4,-1],2)")}
is \itv{1}{16}, and \code{textToInterval("pown([2,5],5)")} throws.
\end{Ex}

\section{What is not provided}\label{sec:ieee1788-missing}
%%----------------------------
\cindex{IEEE 1788-2015!operations not provided}\cindex{decorations}%

GAOL has no decorations: only bare intervals are provided (Clause 11). Nor
does it provide the slope functions (Table~10.6), the reduction operations
\ieee{sum}, \ieee{dot}, \ieee{sumSquare} and \ieee{sumAbs} (12.12.12), the
exact reduction operations of 12.13.5, the conversion specifier \var{cs} of
\ieee{intervalToText} (13.3), or the exceptions \ieee{UndefinedOperation} and
\ieee{PossiblyUndefinedOperation} that \ieee{textToInterval} and
\ieee{exactToInterval} signal (12.12.7, 13.4): for a string that is no
interval they return the empty set silently. Where they cannot tell whether
$l \le u$ in a literal \code{[l, u]}, case (b) of 12.12.7 (the bounds of
\pxref{sec:reading-intervals}), they return the interval from the lower bound
of $l$ to the upper bound of $u$, without \ieee{PossiblyUndefinedOperation}.
That interval holds $[l, u]$ where $l \le u$; where $l$ is above $u$ it is a
false positive that the standard permits, and it can miss $l$ or $u$ when a
bound is computed: \code{"[1/3, 0x1.5555555555555p-2]"} is the point
\code{0x1.5555555555555p-2}, below $1/3$. The end of
\file{gaol/gaol\_ieee1788.h} keeps this list.

\begin{Rem}\label{rem:ieee1788-erf}
In the other direction, a function of GAOL that the standard does not name
is not in \code{gaol\_ieee1788} either, but for \code{sinhRev} and
\code{tanhRev}, which are named after \ieee{coshRev}
(Table~\ref{tab:ieee1788-operations}), for \code{centeredForm(m, r)}, the
constructor \code{interval::centered\_form(m, r)} of a midpoint and a radius
(\pxref{sec:constructors}), and for the operators of intervals, the
relational division \code{\%} among them, which the namespace names as
\code{gaol\_core} does. \code{erf} and \code{erfc}, the
error function and its complement, are two functions of GAOL that the
standard does not name.\newinvfive{} A call
\code{erf(x)} on an interval finds them all the same in \code{gaol\_core},
by argument-dependent lookup, and \code{erf(0.5)} remains the \code{erf}
of C.
\end{Rem}

\section{Discussion}
%%------------------

This chapter has presented the namespace \code{gaol\_ieee1788}, which lets a
program use GAOL under the names and the argument orders of IEEE 1788-2015
with a single directive. The functions that the standard and GAOL name and
define alike are shared between the two namespaces. Where they differ, the
power, the bounds of the empty set, the membership of an infinity, the
writing and the reading of intervals, the namespace follows the standard, and Table~\ref{tab:ieee1788-operations}
gives the translation of every other name. Two operations have their
details here: \code{mulRevToPair}, which \code{gaol} does not name, and the
reverse functions of \code{atan2}, whose arguments \code{gaol\_ieee1788}
orders as \ieee{mulRev} does and \code{gaol} as \code{div\_rel} does.

\medskip
The choice of two namespaces rather than one keeps the programs written for
GAOL~4 and, but for the results that change (Chapter~\ref{chap:changes}),
their results, while a program written for the standard, or compared with
another library of the standard, opens \code{gaol\_ieee1788}. Its price is
stated in \exref{sec:ieee1788-namespace}: a program opens one of the two
namespaces, never both, and qualifies the names of its own that are also names
of the standard. What GAOL does not provide of the standard, the decorations, the
slope functions, the reductions, exact or not, the conversion specifier of
\ieee{intervalToText} and the exceptions of the reading, is listed rather
than approximated (\pxref{sec:ieee1788-missing}), and the header keeps the list.

\medskip
Whatever the name a program calls an operation by, the bounds it gets are
the same. The next chapter states how tight they are: the conditions GAOL
guarantees on each bound, and the tightness of each family of functions.
